\chapter{Conclusiones}
\label{ch:conclusiones}

\section*{Fase~1: Servoregulador}

Los cuatro métodos regulan correctamente la salida $y(k)$ ante escalón y perturbación de carga, pero con diferencias de rendimiento muy marcadas.

El Q-Learning obtiene el menor IAE ($1.80$), con cero sobreimpulso, $t_s=10$~muestras y recuperación instantánea ($0$~ciclos). Su debilidad es el mando: satura ($u_{max}=1.0$) y su variación acumulada ($\sum|\Delta u|=7.10$) es la mayor con diferencia, por el comportamiento bang-bang de las 8~acciones discretas; además, su codificación binaria de la perturbación lo hace frágil ante una rampa no vista (Sección~\ref{sec:robustez_rampa}).

La red profunda, entrenada por optimización directa mediante el método adjunto sobre la planta, ofrece el mejor compromiso: IAE~$=1.90$ ($2\times$ menor que el PID), cero sobreimpulso, $t_s=10$~muestras, el mando más suave de los cuatro ($\sum|\Delta u|=1.65$) y sin saturar ($u_{max}=0.863 < 1.0$). El gradiente del IAE calculado a través de la dinámica de la planta (ecuaciones~\ref{eq:adjoint}--\ref{eq:grad_u}) dirige los pesos hacia la política directamente, y la entrada continua de la perturbación la hace la más robusta ante la rampa.

El controlador difuso, con su base ampliada a 35~reglas, logra el menor consumo energético y un actuador conservador ($u_{max}=0.816$). Es el más adecuado cuando el actuador tiene restricción severa de velocidad de cambio, aunque el más lento en establecerse ($t_s=19$, Rec=12) y el de mayor IAE ($5.97$), al carecer de acción integral.

El PID, tras el reajuste de ganancias, obtiene IAE~$=3.89$, establecimiento en $16$~muestras y recuperación instantánea ($0$~ciclos) por su acción integral. Supera al difuso; las políticas de aprendizaje (DL y QL) siguen por delante en IAE.

\section*{Fase~2: Coordinación R2ET}

Ninguno de los cuatro métodos viola la capacidad máxima de las esteras ni supera el umbral de alarma. La capa de supervisión de capacidad asegura ese resultado con independencia del controlador primario.

El PID-Aware, tras el reajuste de ganancias, obtiene el menor IAE total ($11.41$) y el menor IAE durante el atasco ($1.07$), con el menor desbalance ($0.011$). La integral empuja $n_i$ hacia $n^*$ sin error residual; el supervisor compensa que no anticipe el atasco.

El DL-Aware, con la red más profunda ($8\!\to\!96\!\to\!\dots\!\to\!3$) y 1200~épocas, baja a IAE~$=22.14$ (segundo) y aporta el mayor margen de seguridad ($1.43$~piezas, ocupación máxima de solo $1.57$) y la mayor productividad ($74.93$). El método adjunto de dos esteras converge bien con la capacidad y las épocas adecuadas.

El Q-Learning-Aware combina IAE~$=25.75$ (tercero) con el menor consumo energético ($163.52$) y la mayor robustez ante el escenario severo. Su política opera en la zona baja de ocupación sin gastar velocidad innecesaria; su punto débil es el IAE durante el atasco ($6.12$), por la codificación binaria de la perturbación.

El difuso-Aware, con 35~reglas, reduce su IAE de $56.3$ (versión de 15~reglas) a $37.2$ y mantiene un IAE durante el atasco bajo ($1.82$) y un desbalance mínimo ($0.012$). Al no tener acción integral, queda por detrás de PID, DL y QL en IAE total, pero sus reglas son auditables sin datos ni entrenamiento.

\section*{Supervisor inteligente}

El supervisor de dos capas (ajuste dinámico de ganancias PID y capa dura de seguridad de capacidad) mantiene cero violaciones en todos los métodos ante el cambio de peso. Sin esta capa, ningún controlador escalar coordina el reparto del robot ni reduce la admisión en el momento correcto.

\section*{Recomendación para la CMF R2ET}

Los resultados de ambas fases confirman que la elección del controlador depende del criterio de desempeño prioritario.

Para la capa servo local (regulación de una sola estera, Fase~1), la red profunda con ajuste directo por adjunto es la mejor elección: supera al PID en IAE, no satura, y produce un mando suave. Una vez entrenada, la inferencia es prácticamente instantánea y no requiere resintonización.

Para la coordinación aware R2ET (dos esteras, cambio de peso), la elección depende del criterio prioritario: el PID-Aware retuneado da el menor IAE y es el más simple de mantener; el DL aporta el mayor margen de seguridad y productividad; el Q-Learning es el más eficiente en energía y el más robusto ante cambios severos de condiciones. Para un despliegue auditable sin datos ni entrenamiento, el difuso de 35~reglas es una opción transparente, a costa de un IAE total mayor.

\section*{Contribución metodológica}

La comparación bajo condiciones idénticas (misma planta, escenario y semilla) atribuye las diferencias de rendimiento al controlador y no a condiciones de prueba distintas. En el lazo servo, la red profunda reduce el IAE en un factor de $2\times$ respecto al PID manteniendo el actuador sin saturar y con la menor variación de mando. En la coordinación de dos esteras, dotar a la red de mayor profundidad y más épocas baja su IAE de $75$ a $22$, confirmando que la capacidad y el presupuesto de entrenamiento eran el factor limitante.
